\section{Introduction}
\label{sec:introduction}

Optical recordings can expose the coordinated activity of large neuronal populations, including whole-brain dynamics in larval zebrafish \citep{ahrens2013wholebrain}. Their apparent scale, however, does not remove a basic measurement problem: a fluorescent transient observed at a labeled location may reflect a neuron, nearby neuropil, motion, shared illumination or acquisition effects, or a mixture of these sources. Mature calcium-imaging workflows therefore separate registration, spatial source extraction, background modeling, temporal trace estimation, and event inference rather than treating them as one interchangeable detection problem \citep{pnevmatikakis2016simultaneous,giovannucci2019caiman,zhou2018cnmfe,keemink2018fissa,friedrich2017oasis}.

The present study occupies a data regime in which conventional supervised-learning claims are especially difficult to support. The available canonical recording contains a small number of shared burst intervals, and the positive annotations identify only a subset of plausible neurons. The labels are useful anchors but do not define an exhaustive negative class. Consequently, an unmatched algorithmic candidate cannot automatically be called a false positive, and the absence of a label cannot be interpreted as biological inactivity. This is a positive--unlabeled setting in which the labeling mechanism itself may favor bright, isolated, or otherwise readily visible neurons \citep{bekker2020pu}.

These constraints motivate an identifiability-first question. Rather than asking only whether an operator assigns a high score at a labeled location, we ask which components of the observed signal are temporally exceptional, spatially localized, reproducible within a site, shared across the field, and stable across repeated bursts. This distinction is essential because a representation can be useful for detecting change while discarding fluorescence amplitude, distorting peak time, or retaining the same waveform in surrounding tissue. Likewise, a high event-aligned correlation can reflect a recording-level burst template without establishing a neuron-specific activation signature.

The repository already contains evidence that illustrates this distinction. A learned two-frame independent component is nearly collinear with signed temporal difference; event-aligned Raw traces contain a strong shared waveform; and local background subtraction increases event-associated residual amplitude without yet establishing a universal residual morphology. These results are informative boundary conditions, but the current identity handling, sparse labels, and reuse of the same four bursts require a protected confirmatory analysis before final numerical claims are made.

We therefore develop a reproducible analysis program with four objectives. First, we separate immutable observation-site geometry from provisional canonical-neuron identity and export a complete Raw-to-pipeline trace atlas. Second, we characterize acquisition noise, common-mode activity, and spatial contamination before interpreting neuron-level variation. Third, we partition scalar and functional variability into site-, burst-, and residual components and test whether amplitude, signal-to-noise ratio, timing, and spatial compactness define stable measurement phenotypes. Fourth, we evaluate unsupervised representations using positive--unlabeled-compatible metrics and reserve precision estimation for a bounded field that receives exhaustive blinded review. The intended contribution is not a claim that a single operator solves neuron detection, but a framework for determining what the available recording can and cannot identify.
